% Chapter 1, typeset from lectures/L01: the README is the opener and the closing summary, and
% each appendix is a section, in the same order and with the same subsections. Appendix D, the
% solutions, is in back/answers/c01.tex.

\chapter{Talsystem och binära koder}
\label{c1:ch}

{\large\itshape Varför en dator räknar med ettor och nollor, och hur man räknar som den.\par}

\begin{chapterbox}{I det här kapitlet}
\begin{itemize}
  \item Digitalt och analogt: varför nästan all teknik i dag räknar med ettor och nollor.
  \item Positionssystem: det decimala, det binära och det hexadecimala talsystemet.
  \item Omvandling mellan decimalt, binärt och hexadecimalt, för hand.
  \item Bitar, byte och talområden: hur stora tal som ryms i 4, 8 och 16 bitar.
  \item Binär addition, och vad som händer när summan inte ryms.
  \item Binära koder: BCD, Graykod, ASCII, 7-segmentkod och paritetsbit.
\end{itemize}
\end{chapterbox}

Allt i den här boken vilar på ett enda faktum: en digital signal har bara två tillstånd. Det här
kapitlet visar varför det är en styrka och inte en begränsning, och hur många sådana signaler
tillsammans kan bära ett tal, en bokstav eller en siffra på en display.

Arbeta igenom kapitlet i ordning. \Secref{c1:app:a} om talsystemen är grunden för resten av boken:
viktraden 1, 2, 4, 8, 16, 32, 64, 128 och omvandlingarna mellan decimalt, binärt och hexadecimalt
används i varje kapitel som följer, och de sitter först när du har räknat många själv.
\Secref{c1:app:b} går igenom de binära koder du kommer att möta, både i kursen och i yrkeslivet.
Övningarna i \secref{c1:sec:exercises} räknas med papper och penna, och deras lösningsförslag
finns i bilaga~\ref{app:answers}.

% Section 1.1, from lectures/L01/appendix/a_number_systems.md.

\section{Talsystem}
\label{c1:app:a}

\subsection{Digitalt och analogt}
\label{c1:a1}
En \textbf{analog} signal kan anta vilket värde som helst inom ett intervall. Spänningen från en
temperaturgivare kan vara 1,2~V, 1,21~V eller 1,2137~V, och varje liten ändring betyder något. Det
gör också varje liten störning: brus från en motor i närheten flyttar signalen, och mottagaren
kan inte veta vilken del av värdet som är mätning och vilken som är störning.

En \textbf{digital} signal antar bara två värden, som vi kallar \code{0} och \code{1}. I en krets
som matas med 5~V betyder ungefär 0~V en nolla och ungefär 5~V en etta, och allt som ligger
tillräckligt nära någon av dem läses som exakt \code{0} eller exakt \code{1}. En störning på några
tiondels volt ändrar därför ingenting alls.

\textbf{Den tåligheten mot brus är skälet till att nästan all modern elektronik är digital}, från
mikrokontrollern i en kaffebryggare till styrsystemet i en fabrik. Priset är att en digital signal
bara bär en enda uppgift, ja eller nej, tänd eller släckt. För att kunna representera ett tal, en
temperatur eller en bokstav behövs många digitala signaler tillsammans, och ett sätt att tolka
dem. Det är vad talsystemen och koderna i det här kapitlet handlar om.

En digital signal, eller en siffra i ett binärt tal, kallas en \textbf{bit}, av engelskans
\emph{binary digit}.

\subsection{Positionssystem och talbas}
\label{c1:a2}
Det talsystem vi använder till vardags, det \textbf{decimala} talsystemet, har tio siffror, 0
till 9. Antalet siffror kallas talsystemets \textbf{talbas}, så det decimala talsystemet har
talbasen 10.

Det är också ett \textbf{positionssystem}: en siffras värde beror på var i talet den står. I talet
25 är femman värd 5 och tvåan värd 20. Lägg till en nolla sist, 250, så flyttar båda ett steg åt
vänster och blir värda tio gånger mer: femman 50 och tvåan 200.

Varje position har alltså en \textbf{vikt}, en potens av talbasen. Positionerna numreras från
höger med början på \textbf{noll}:

\begin{narrowtable}{@{}lccc@{}}
  \thead{Position} & 2 & 1 & 0 \\
  \midrule
  Vikt & $10^2 = 100$ & $10^1 = 10$ & $10^0 = 1$ \\
  Siffra i talet 250 & 2 & 5 & 0 \\
  Värde & 200 & 50 & 0 \\
\end{narrowtable}

\noindent Talets värde är summan av varje siffra gånger vikten för dess position:
\[
  250 = 2 \cdot 10^2 + 5 \cdot 10^1 + 0 \cdot 10^0 = 200 + 50 + 0
\]

Att numreringen börjar på noll är ingen godtycklig vana. Den högra positionen har vikten
$B^0 = 1$ i \textbf{alla} talsystem, oavsett talbas $B$, och det är skälet till att den kallas
position noll.

Samma formel gäller alltså för varje talsystem. En siffra $x_n$ på position $n$ i ett tal med
talbasen $B$ är värd
\[
  x_n \cdot B^n
\]
och talets värde är summan av alla siffrornas värden. Den högra siffran, med lägst vikt, kallas
den \textbf{minst signifikanta} siffran, och den vänstra den \textbf{mest signifikanta}.

När det inte framgår av sammanhanget vilket talsystem ett tal är skrivet i, skrivs talbasen som
ett index efter talet: $25_{10}$ är tjugofem decimalt, $11001_2$ är samma tal binärt. I kod och i
den här bokens text används i stället ett prefix, se \secref{c1:a8}.

\subsection{Det binära talsystemet}
\label{c1:a3}
En digital signal har två tillstånd, så den naturliga talbasen i digitaltekniken är 2. Det
\textbf{binära} talsystemet har bara siffrorna 0 och 1, och varje position är värd dubbelt så
mycket som positionen till höger om den:

\begin{narrowtable}{@{}lcccccccc@{}}
  \thead{Position} & 7 & 6 & 5 & 4 & 3 & 2 & 1 & 0 \\
  \midrule
  Vikt & 128 & 64 & 32 & 16 & 8 & 4 & 2 & 1 \\
\end{narrowtable}

\noindent Den raden med vikter är det viktigaste i hela det här kapitlet. Lär dig den utantill,
från höger: 1, 2, 4, 8, 16, 32, 64, 128.

\subsubsection{Från binärt till decimalt}
För att omvandla ett binärt tal till decimalt summerar du vikterna för de positioner där det står
en etta. Nollorna bidrar inte med något och kan hoppas över.

\textbf{Exempel:} $1011_2$ har ettor på position 3, 1 och 0:
\[
  1011_2 = 1 \cdot 8 + 0 \cdot 4 + 1 \cdot 2 + 1 \cdot 1 = 8 + 2 + 1 = 11_{10}
\]

\textbf{Exempel:} $1010\,1110_2$ har ettor på position 7, 5, 3, 2 och 1:

\bookfigure{l01_bit_weights}{Ett 8-bitars tal med bitnumren ovanför och vikterna under: talet
\code{1010 1110} summeras till 174.}{c1:fig:bit_weights}

\[
  1010\,1110_2 = 128 + 32 + 8 + 4 + 2 = 174_{10}
\]

Binära tal skrivs ofta i grupper om fyra siffror, med ett mellanrum emellan, på samma sätt som man
skriver 1\,000\,000 i stället för 1000000. Mellanrummet betyder ingenting; det gör bara talet
lättare att läsa.

\subsubsection{Bitar, nibblar och byte}
Några namn på grupper av bitar som återkommer genom hela boken:

\begin{narrowtable}{@{}lcl@{}}
  \thead{Namn} & \thead{Antal bitar} & \thead{Exempel} \\
  \midrule
  bit & 1 & \code{1} \\
  nibble & 4 & \code{1011} \\
  byte & 8 & \code{1010 1110} \\
  ord (på AVR-processorn) & 16 & \code{0000 0011 1110 1000} \\
\end{narrowtable}

\noindent Den vänstra biten, med störst vikt, kallas \textbf{MSB} (\emph{most significant bit}).
Den högra, med vikten 1, kallas \textbf{LSB} (\emph{least significant bit}). I en byte är MSB
bit~7 och LSB bit~0.

Två saker går att se direkt på LSB och MSB:
\begin{itemize}
  \item \textbf{LSB avgör om talet är udda eller jämnt.} Alla andra vikter är jämna tal, så bara
    LSB kan göra summan udda.
  \item \textbf{MSB avgör om talet är minst hälften av det största talet.} I en byte är MSB värd
    128, och resten tillsammans som mest 127.
\end{itemize}

\subsubsection{Talområden}
Med $n$ bitar finns det $2^n$ olika kombinationer av ettor och nollor, eftersom varje bit
fördubblar antalet möjligheter. Den minsta är noll, så den största är en mindre än antalet
kombinationer:
\[
  \text{antal värden} = 2^n \qquad \text{största värde} = 2^n - 1
\]

\begin{narrowtable}{@{}ccc@{}}
  \thead{Antal bitar} & \thead{Antal värden} & \thead{Talområde} \\
  \midrule
  1 & 2 & 0--1 \\
  4 & 16 & 0--15 \\
  8 & 256 & 0--255 \\
  10 & 1024 & 0--1023 \\
  16 & 65\,536 & 0--65\,535 \\
\end{narrowtable}

\noindent Du kan kontrollera den största: åtta ettor är $128 + 64 + 32 + 16 + 8 + 4 + 2 + 1 =
255$, vilket är $2^8 - 1$.

Det här är ingen teoretisk detalj. Mikrodatorn i den här boken, ATmega328P, räknar med 8 bitar åt
gången, så ett av dess register kan bara hålla talen 0--255. Vad som händer när ett tal inte ryms
är ämnet för \secref{c1:a7}, och för en stor del av kapitel~\ref{c8:ch}.

\subsection{Från decimalt till binärt}
\label{c1:a4}
Åt andra hållet finns två metoder. Båda ger samma svar; välj den du tycker är enklast.

\subsubsection{Metod 1: vikttabellen}
Skriv upp vikterna, och ta för varje vikt från vänster ställning till om den ryms i det som är kvar
av talet. Ryms den skriver du en etta och drar bort vikten; annars skriver du en nolla.

\textbf{Exempel:} omvandla $23_{10}$ till binärt.

\begin{narrowtable}{@{}lccccc@{}}
  \thead{Vikt} & 16 & 8 & 4 & 2 & 1 \\
  \midrule
  Ryms den i det som är kvar? & 23: ja & 7: nej & 7: ja & 3: ja & 1: ja \\
  Kvar efteråt & 7 & 7 & 3 & 1 & 0 \\
  Bit & 1 & 0 & 1 & 1 & 1 \\
\end{narrowtable}

\noindent Alltså är $23_{10} = 10111_2$. Kontrollera baklänges: $16 + 4 + 2 + 1 = 23$.

\subsubsection{Metod 2: division med 2}
Dela talet med 2 och skriv upp resten, som alltid är 0 eller 1. Dela sedan kvoten med 2 igen, och
fortsätt tills kvoten är 0. Resterna är bitarna, med \textbf{den första resten som LSB}.

\textbf{Exempel:} omvandla $23_{10}$ till binärt igen.

\begin{narrowtable}{@{}ccl@{}}
  \thead{Division} & \thead{Kvot} & \thead{Rest} \\
  \midrule
  23 / 2 & 11 & 1 (bit 0, LSB) \\
  11 / 2 & 5 & 1 (bit 1) \\
  5 / 2 & 2 & 1 (bit 2) \\
  2 / 2 & 1 & 0 (bit 3) \\
  1 / 2 & 0 & 1 (bit 4, MSB) \\
\end{narrowtable}

\noindent Resterna läses \textbf{nedifrån och upp}: $10111_2$, samma svar som med tabellen.

Det vanligaste felet med den här metoden är att läsa resterna uppifrån och ned, vilket ger talet
baklänges, $11101_2 = 29$. Kontrollera därför alltid svaret genom att räkna tillbaka till
decimalt.

\subsubsection{Att fylla ut till en hel byte}
$10111_2$ har fem bitar. Ska talet ligga i ett 8-bitars register fylls det ut med nollor till
vänster: \code{0001 0111}. Nollor till vänster ändrar inte värdet, precis som 023 är samma sak som
23.

\subsection{Det hexadecimala talsystemet}
\label{c1:a5}
Binära tal blir snabbt långa och svårlästa. Ett 16-bitars tal som \code{0100 1010 1000 1100} är
lätt att skriva av fel, och det är nästan omöjligt att se vilket tal det är. Därför skrivs binära
tal nästan alltid i det \textbf{hexadecimala} talsystemet, med talbasen 16.

Det behövs 16 siffror, så efter 0--9 används bokstäverna A--F:

\begin{narrowtable}{@{}ccc@{\hspace{3em}}ccc@{}}
  \thead{Decimalt} & \thead{Binärt} & \thead{Hexadecimalt} &
  \thead{Decimalt} & \thead{Binärt} & \thead{Hexadecimalt} \\
  \midrule
  0 & \code{0000} & 0 & 8 & \code{1000} & 8 \\
  1 & \code{0001} & 1 & 9 & \code{1001} & 9 \\
  2 & \code{0010} & 2 & 10 & \code{1010} & A \\
  3 & \code{0011} & 3 & 11 & \code{1011} & B \\
  4 & \code{0100} & 4 & 12 & \code{1100} & C \\
  5 & \code{0101} & 5 & 13 & \code{1101} & D \\
  6 & \code{0110} & 6 & 14 & \code{1110} & E \\
  7 & \code{0111} & 7 & 15 & \code{1111} & F \\
\end{narrowtable}

\noindent\textbf{Poängen med just talbasen 16 är att $16 = 2^4$: en hexadecimal siffra motsvarar
exakt fyra bitar}, en nibble. Omvandlingen mellan binärt och hexadecimalt kräver därför ingen
räkning alls, bara tabellen ovan.

\subsubsection{Från binärt till hexadecimalt}
Dela upp det binära talet i grupper om fyra bitar, \textbf{med början från höger}, och byt varje
grupp mot sin hexadecimala siffra.

\textbf{Exempel:} $1011\,0110_2$.

\bookfigure{l01_hex_grouping}{Talet \code{1011 0110} uppdelat i två grupper om fyra bitar: den
vänstra gruppen, \code{1011}, blir B och den högra, \code{0110}, blir 6.}{c1:fig:hex_grouping}

\noindent Alltså är $1011\,0110_2 = \text{B6}_{16}$.

Om den sista gruppen till vänster har färre än fyra bitar fyller du ut den med nollor:
$1\,1110_2 = 0001\,1110_2 = \text{1E}_{16}$. Det är därför grupperingen börjar från höger. Börjar
man från vänster hamnar de utfyllande nollorna mitt i talet, där de ändrar värdet.

\subsubsection{Från hexadecimalt till binärt}
Byt varje hexadecimal siffra mot sina fyra bitar.

\textbf{Exempel:} $\text{4A8C}_{16}$.

\begin{narrowtable}{@{}lcccc@{}}
  \thead{Hexadecimal siffra} & 4 & A & 8 & C \\
  \midrule
  Fyra bitar & \code{0100} & \code{1010} & \code{1000} & \code{1100} \\
\end{narrowtable}

\noindent Alltså är $\text{4A8C}_{16} = 0100\,1010\,1000\,1100_2$. Fyra tecken i stället för
sexton, och omvandlingen tar några sekunder när tabellen väl sitter.

\subsection{Omvandling mellan alla tre}
\label{c1:a6}
Mellan binärt och hexadecimalt är det bara att gruppera, som i \secref{c1:a5}. Mellan decimalt och
de andra två finns två vägar.

\subsubsection{Från hexadecimalt till decimalt}
Använd formeln från \secref{c1:a2} med talbasen 16. Vikterna är $16^0 = 1$, $16^1 = 16$,
$16^2 = 256$ och $16^3 = 4096$.

\textbf{Exempel:} $\text{F6}_{16}$, där F är 15:
\[
  \text{F6}_{16} = 15 \cdot 16 + 6 \cdot 1 = 240 + 6 = 246_{10}
\]

\textbf{Exempel:} $\text{2A5}_{16}$, där A är 10:
\[
  \text{2A5}_{16} = 2 \cdot 256 + 10 \cdot 16 + 5 \cdot 1 = 512 + 160 + 5 = 677_{10}
\]

Den andra vägen går via binärt: $\text{F6}_{16} = 1111\,0110_2 = 128 + 64 + 32 + 16 + 4 + 2 =
246$. Den är ofta lika snabb för tvåsiffriga tal, och bra som kontroll.

\subsubsection{Från decimalt till hexadecimalt}
Också här finns två vägar.

\textbf{Via binärt:} omvandla till binärt enligt \secref{c1:a4}, och gruppera sedan i fyror.
$200_{10}$ blir $1100\,1000_2$, som är $\text{C8}_{16}$.

\textbf{Division med 16:} som divisionsmetoden i \secref{c1:a4}, men med 16 i stället för 2.
Resterna är hexadecimala siffror, och även här är den första resten den minst signifikanta.

\begin{narrowtable}{@{}ccl@{}}
  \thead{Division} & \thead{Kvot} & \thead{Rest} \\
  \midrule
  200 / 16 & 12 & 8 \\
  12 / 16 & 0 & 12 = C \\
\end{narrowtable}

\noindent Nedifrån och upp: $\text{C8}_{16}$. Kontroll: $12 \cdot 16 + 8 = 192 + 8 = 200$.

\subsubsection{En sammanfattning}

\begin{booktable}{@{}p{7.2em}LLL@{}}
  \thead{Från \textbackslash{} Till} & \thead{Decimalt} & \thead{Binärt} & \thead{Hexadecimalt} \\
  \midrule
  \textbf{Decimalt} & -- & vikttabellen eller division med 2 & via binärt, eller division med
    16 \\
  \textbf{Binärt} & summera vikterna & -- & gruppera i fyror från höger \\
  \textbf{Hexadecimalt} & siffra gånger $16^n$ & fyra bitar per siffra & -- \\
\end{booktable}

\subsection{Binär addition}
\label{c1:a7}
Binära tal adderas med samma uppställning som decimala: kolumn för kolumn från höger, med en
\textbf{minnessiffra} (en \emph{carry}) som förs vidare till nästa kolumn när summan inte ryms i
en siffra.

Skillnaden är bara att det \enquote{blir tio} redan vid två. Fyra regler räcker:

\begin{narrowtable}{@{}lcc@{}}
  \thead{Summa} & \thead{Resultat} & \thead{Minnessiffra} \\
  \midrule
  0 + 0 & 0 & 0 \\
  0 + 1 & 1 & 0 \\
  1 + 1 & 0 & 1 \\
  1 + 1 + 1 (med minnessiffra) & 1 & 1 \\
\end{narrowtable}

\noindent Den tredje raden är den som känns ovan: $1 + 1 = 10_2$, alltså två, skrivet som en nolla
med en etta i minnet.

\textbf{Exempel:} $90 + 55$, med båda talen som 8-bitars binära tal.

\begin{textdrawing}
  minne:   1111 1100
           0101 1010     (90)
         + 0011 0111     (55)
         -----------
           1001 0001     (145)
\end{textdrawing}

\noindent Kontrollera: $1001\,0001_2 = 128 + 16 + 1 = 145$, och $90 + 55 = 145$.

\subsubsection{När summan inte ryms}
Ett 8-bitars register rymmer 0--255. Vad händer med $240 + 17 = 257$?

\begin{textdrawing}
  minne: 1 1110 0000
           1111 0000     (240)
         + 0001 0001     (17)
         -----------
         1 0000 0001
\end{textdrawing}

\noindent Summan behöver nio bitar. I ett 8-bitars register finns bara plats för de åtta högra,
\code{0000 0001}, alltså 1, och den nionde biten, minnessiffran ut ur bit~7, hamnar utanför.
Resultatet blir $257 - 256 = 1$.

\textbf{Det är inte ett fel i datorn; det är vad åtta bitar betyder.} Talet slår runt, som
vägmätaren i en gammal bil som går från 99\,999 till 00\,000. Mikrodatorn sparar den nionde biten
i en särskild flagga, \emph{carry}, så att programmet kan upptäcka att det hände. Det är ämnet för
kapitel~\ref{c8:ch}, där samma sak kallas \textbf{aritmetisk rundgång}.

\subsubsection{Hexadecimal addition}
Hexadecimala tal kan adderas direkt, med regeln att det \enquote{blir tio} vid 16. Oftast är det
lika enkelt att gå via binärt eller decimalt.

\textbf{Exempel:} $\text{3C}_{16} + \text{0A}_{16}$. Entalskolumnen: C + A = 12 + 10 = 22 = 16 + 6,
så siffran blir 6 med en etta i minnet. Sextontalskolumnen: 3 + 0 + 1 = 4. Summan är
$\text{46}_{16}$, och kontrollen i decimalt är $60 + 10 = 70 = 4 \cdot 16 + 6$.

\subsection{Skrivsätt}
\label{c1:a8}
Samma tal kan skrivas på flera sätt, och vilket som används beror på sammanhanget.

\begin{booktable}{@{}p{8.6em}p{10.4em}L@{}}
  \thead{Skrivsätt} & \thead{Exempel (tjugofem)} & \thead{Används} \\
  \midrule
  Index efter talet & $25_{10}$, $11001_2$, $19_{16}$ & i matematik och i läroböcker \\
  Prefix \code{0b} & \code{0b00011001} & i kod, för binära tal \\
  Prefix \code{0x} & \code{0x19} & i kod, för hexadecimala tal \\
  Prefix \code{$} & \code{$19} & i äldre AVR-assembler, för hexadecimala tal \\
  Utan prefix & \code{25} & alltid decimalt \\
\end{booktable}

Den här boken använder prefixen \code{0b} och \code{0x} i löptext och i kod, eftersom det är så
talen skrivs i assemblerprogrammen från kapitel~\ref{c7:ch} och framåt. Indexen används bara i
härledningar som den här, där båda skrivsätten behövs sida vid sida.

Två saker är värda att notera:
\begin{itemize}
  \item \code{0x19} och \code{19} är olika tal: tjugofem respektive nitton. Ett tal utan prefix är
    decimalt.
  \item Nollor till vänster ändrar inte värdet: \code{0x09} och \code{0x9} är samma tal. I kod
    skriver man ofta ut dem ändå, så att det syns hur många bitar talet är tänkt att ha.
\end{itemize}

\subsubsection{Att kontrollera med kalkylatorn}
Kalkylatorn i Windows har ett läge som heter \textbf{Programmerare}. Där visas samma tal samtidigt
som \code{HEX}, \code{DEC}, \code{OCT} (oktalt, talbas 8, som inte används i boken) och \code{BIN}.
Knappen där det står \code{QWORD}, under talen, väljer hur många bitar kalkylatorn räknar med;
varje klick byter till nästa: \code{QWORD} (64), \code{DWORD} (32), \code{WORD} (16) och
\code{BYTE} (8).

Använd den för att kontrollera dina omvandlingar, aldrig för att ta fram dem: på provet finns ingen
kalkylator, och en omvandling du inte kan göra för hand kan du inte heller felsöka i ett program.

% Section 1.2, from lectures/L01/appendix/b_binary_codes.md.

\section{Binära koder}
\label{c1:app:b}

\subsection{Vad en kod är}
\label{c1:b1}
Ett binärt tal är ett sätt att tolka en rad bitar: som en summa av vikter. Men samma rad bitar kan
tolkas på andra sätt, och vilket sätt som gäller är en överenskommelse mellan den som skriver och
den som läser. En sådan överenskommelse kallas en \textbf{binär kod}.

Bitmönstret \code{0100 0001} är talet 65 om det tolkas som ett binärt tal, bokstaven \code{A} om
det tolkas som ASCII, och talet 41 om det tolkas som BCD. Bitarna vet inte vilket de är.
\textbf{Det är alltid sammanhanget, och den som skrev programmet, som bestämmer vad ett bitmönster
betyder.}

Det här avsnittet går igenom de koder som förekommer i boken och som du kommer att möta i
yrkeslivet:

\begin{booktable}{@{}p{6.4em}LL@{}}
  \thead{Kod} & \thead{Kodar} & \thead{Används till exempel i} \\
  \midrule
  BCD & decimala siffror & displayer, klockkretsar, äldre styrsystem \\
  Graykod & tal, så att bara en bit ändras åt gången & lägesgivare, Karnaughdiagram
    (kapitel~\ref{c4:ch}) \\
  ASCII & bokstäver, siffertecken och skiljetecken & all text i datorer, seriell kommunikation \\
  7-segmentkod & vilka segment som ska lysa i en sifferdisplay & displayer på instrument och
    apparater \\
  Paritetsbit & en extra bit som avslöjar ett överföringsfel & seriell kommunikation, minnen \\
\end{booktable}

\subsection{BCD}
\label{c1:b2}
\textbf{BCD}, \emph{binary coded decimal}, kodar varje decimal siffra för sig med fyra bitar, med
vikterna 8, 4, 2 och 1. Koden kallas därför också \textbf{8421-kod}.

\begin{narrowtable}{@{}l*{10}{c}@{}}
  \thead{Decimal siffra} & 0 & 1 & 2 & 3 & 4 & 5 & 6 & 7 & 8 & 9 \\
  \midrule
  BCD & \code{0000} & \code{0001} & \code{0010} & \code{0011} & \code{0100} & \code{0101} &
    \code{0110} & \code{0111} & \code{1000} & \code{1001} \\
\end{narrowtable}

\noindent Ett tal med flera siffror kodas siffra för siffra, fyra bitar per siffra:

\begin{console}
59 (decimalt)  ->  5 = 0101, 9 = 1001  ->  0101 1001 (BCD)
\end{console}

\noindent Jämför med det binära talet 59, som är \code{0011 1011}. \textbf{BCD och binärt är olika
koder}, och förväxlar man dem blir det fel: \code{0101 1001} tolkat som ett binärt tal är 89.

Fyra bitar har 16 kombinationer, men BCD använder bara tio av dem. Kombinationerna \code{1010}
till \code{1111} är \textbf{ogiltiga} i BCD. Det är slöseri med bitar, men priset är värt det när
talet ska visas: varje nibble är redan en decimal siffra, redo att skickas till sin egen display,
utan någon omräkning. Därför används BCD i klockkretsar och sifferdisplayer. Att koden lämnar sex
kombinationer oanvända återkommer i kapitel~\ref{c4:ch}, där de blir \emph{don't care} i ett
Karnaughdiagram.

\subsection{Graykod}
\label{c1:b3}
Tänk dig en givare som mäter vinkeln på en axel med tre bitar, till exempel i en robotarm. När
axeln vrids från läge 3 till läge 4 ska givarens utsignal gå från \code{011} till \code{100}. Men
alla tre bitarna ändras, och de ändras aldrig på exakt samma gång. Under det korta ögonblick då
bara några av dem har hunnit ändras läser styrsystemet något helt annat, kanske \code{111}, alltså
läge 7. Axeln har inte rört sig mer än en bråkdel av ett steg, men systemet tror att den har gått
halvvägs runt.

\textbf{Graykoden} löser det genom att ordna kombinationerna så att \textbf{exakt en bit ändras
mellan två på varandra följande tal}, också från det sista talet tillbaka till det första:

\begin{narrowtable}{@{}ccc@{}}
  \thead{Tal} & \thead{Binärt} & \thead{Graykod} \\
  \midrule
  0 & \code{000} & \code{000} \\
  1 & \code{001} & \code{001} \\
  2 & \code{010} & \code{011} \\
  3 & \code{011} & \code{010} \\
  4 & \code{100} & \code{110} \\
  5 & \code{101} & \code{111} \\
  6 & \code{110} & \code{101} \\
  7 & \code{111} & \code{100} \\
\end{narrowtable}

\noindent Mellan läge 3 (\code{010}) och läge 4 (\code{110}) ändras nu bara den vänstra biten.
Under övergången kan givaren alltså bara visa 3 eller 4, aldrig något orimligt.

\subsubsection{Att bygga en Graykod}
En Graykod med en bit fler byggs genom att \textbf{spegla} den man har:
\begin{enumerate}
  \item Skriv koden med två bitar: \code{00}, \code{01}, \code{11}, \code{10}.
  \item Skriv samma lista en gång till, i omvänd ordning: \code{10}, \code{11}, \code{01},
    \code{00}.
  \item Sätt en nolla framför den första halvan och en etta framför den andra.
\end{enumerate}
Resultatet är tabellen ovan. Därför kallas koden också \textbf{reflekterad binärkod}.

Tvåbitarsversionen, \code{00}, \code{01}, \code{11}, \code{10}, är den du kommer att använda mest.
Den är ordningen på rader och kolumner i ett Karnaughdiagram, där just egenskapen att grannar
skiljer sig i en enda bit är hela poängen (kapitel~\ref{c4:ch}).

\subsection{ASCII}
\label{c1:b4}
En dator lagrar text som tal: varje tecken har en kod. Den kod som nästan all text bygger på heter
\textbf{ASCII}, \emph{American Standard Code for Information Interchange}. Den använder 7 bitar och
har alltså 128 tecken: bokstäver, siffertecken, skiljetecken och några styrtecken.

Ett urval:

\begin{narrowtable}{@{}lcc@{\hspace{3em}}lcc@{}}
  \thead{Tecken} & \thead{Hexadecimalt} & \thead{Decimalt} &
  \thead{Tecken} & \thead{Hexadecimalt} & \thead{Decimalt} \\
  \midrule
  mellanslag & \code{0x20} & 32 & \code{A} & \code{0x41} & 65 \\
  \code{0} & \code{0x30} & 48 & \code{B} & \code{0x42} & 66 \\
  \code{1} & \code{0x31} & 49 & \code{Z} & \code{0x5A} & 90 \\
  \code{9} & \code{0x39} & 57 & \code{a} & \code{0x61} & 97 \\
  \code{:} & \code{0x3A} & 58 & \code{b} & \code{0x62} & 98 \\
  \code{?} & \code{0x3F} & 63 & \code{z} & \code{0x7A} & 122 \\
\end{narrowtable}

\noindent Styrtecknen ligger på \code{0x00}--\code{0x1F} och skrivs inte ut. De två du oftast möter
är radmatning (\code{0x0A}) och vagnretur (\code{0x0D}), som avslutar en rad text.

Tre saker i tabellen är värda att lägga på minnet:
\begin{itemize}
  \item \textbf{Siffertecknen \code{0}--\code{9} ligger på \code{0x30}--\code{0x39}.} Siffertecknet
    \code{7} är alltså inte talet 7 utan talet \code{0x37} = 55. För att få talet ur tecknet drar
    man bort \code{0x30}, och för att få tecknet ur talet lägger man till \code{0x30}. Det är precis
    vad ett program gör när det ska visa ett tal som text, och det gör du själv i
    kapitel~\ref{c10:ch}.
  \item \textbf{Bokstäverna ligger i ordning}, så \code{B} är ett mer än \code{A}. Det gör att ett
    program kan räkna fram bokstäver.
  \item \textbf{Stora och små bokstäver skiljer sig med \code{0x20}}, alltså bara i bit~5: \code{A}
    är \code{0100 0001} och \code{a} är \code{0110 0001}.
\end{itemize}

ASCII har inga svenska bokstäver. Å, ä och ö kom med senare koder som utökar ASCII, i dag nästan
alltid Unicode i formatet UTF-8, där de första 128 tecknen är exakt ASCII.

\subsection{7-segmentkod}
\label{c1:b5}
En sifferdisplay på en mikrovågsugn, en multimeter eller en klocka består ofta av sju avlånga
lysdioder, \textbf{segment}, som kallas a till g. Varje siffra visas genom att rätt segment tänds.

\bookfigure{l01_seven_segment}{En 7-segmentsdisplay med segmenten a till g, och samma display när
segmenten a, c, d, f och g lyser och bildar siffran 5.}{c1:fig:seven_segment}

\textbf{7-segmentkoden} för en siffra är en bit per segment: \code{1} för tänt och \code{0} för
släckt. Skrivs bitarna i ordningen g, f, e, d, c, b, a, med a som bit~0, blir koderna för siffrorna
0--9:

\begin{narrowtable}{@{}clcc@{}}
  \thead{Siffra} & \thead{Tända segment} & \thead{\code{g f e d c b a}} & \thead{Hexadecimalt} \\
  \midrule
  0 & a b c d e f & \code{0 1 1 1 1 1 1} & \code{0x3F} \\
  1 & b c & \code{0 0 0 0 1 1 0} & \code{0x06} \\
  2 & a b d e g & \code{1 0 1 1 0 1 1} & \code{0x5B} \\
  3 & a b c d g & \code{1 0 0 1 1 1 1} & \code{0x4F} \\
  4 & b c f g & \code{1 1 0 0 1 1 0} & \code{0x66} \\
  5 & a c d f g & \code{1 1 0 1 1 0 1} & \code{0x6D} \\
  6 & a c d e f g & \code{1 1 1 1 1 0 1} & \code{0x7D} \\
  7 & a b c & \code{0 0 0 0 1 1 1} & \code{0x07} \\
  8 & a b c d e f g & \code{1 1 1 1 1 1 1} & \code{0x7F} \\
  9 & a b c d f g & \code{1 1 0 1 1 1 1} & \code{0x6F} \\
\end{narrowtable}

\noindent En krets som tar emot en BCD-siffra och tänder rätt segment kallas en \textbf{avkodare},
och den är ett kombinatoriskt nät: sju utgångar, en per segment, som var och en är en funktion av
de fyra ingångsbitarna. Sådana nät är ämnet för kapitel~\ref{c2:ch} och kapitel~\ref{c4:ch}.

Tabellen gäller en display där en etta tänder segmentet (gemensam katod). Det finns också
displayer där en nolla tänder segmentet (gemensam anod), och då är varje bit inverterad. Vilken sort
man har står i databladet.

\subsection{Paritetsbit}
\label{c1:b6}
När bitar skickas över en kabel kan en störning ibland vända en etta till en nolla eller tvärtom.
Det enklaste sättet att upptäcka det är en \textbf{paritetsbit}: en extra bit som läggs till så att
det totala antalet ettor alltid är jämnt.

\textbf{Exempel:} tecknet \code{A} är \code{100 0001} i ASCII. Det har två ettor, ett jämnt antal,
så paritetsbiten blir \code{0}, och det som skickas är \code{0100 0001}. Tecknet \code{C} är
\code{100 0011} med tre ettor, så paritetsbiten blir \code{1}, och det som skickas är
\code{1100 0011}, med fyra ettor.

Mottagaren räknar ettorna. Är antalet udda har något gått fel, och tecknet kan begäras på nytt.

Paritetsbiten har en tydlig begränsning: \textbf{den upptäcker en felaktig bit, men inte två.}
Vänds två bitar är antalet ettor jämnt igen, och felet passerar obemärkt. Den säger inte heller
vilken bit som är fel. Den är ändå vanlig, eftersom den kostar en enda bit och fångar det
vanligaste felet.

Att räkna ut en paritetsbit är samma sak som att använda XOR-grinden från kapitel~\ref{c2:ch} på
alla bitar: XOR av ett jämnt antal ettor är \code{0}, och av ett udda antal \code{1}.

Det finns också \textbf{udda paritet}, där paritetsbiten väljs så att antalet ettor blir udda.
Sändare och mottagare måste vara överens om vilken som gäller; det är ännu en överenskommelse av det
slag som \secref{c1:b1} handlade om.

% Section 1.3, from lectures/L01/README.md: what you should be able to do, the questions to test
% yourself with, and the next lecture.

\section{Sammanfattning}
\label{c1:sec:summary}

\needspace{6\baselineskip}
\subsection*{Det här ska du kunna nu}
Efter det här kapitlet ska du kunna:
\begin{itemize}
  \item Förklara skillnaden mellan en digital och en analog signal, och varför digitala signaler
    tål brus.
  \item Beskriva det decimala, det binära och det hexadecimala talsystemet: talbas, siffror och
    positionernas vikter.
  \item Omvandla tal mellan decimalt, binärt och hexadecimalt för hand, i alla riktningar.
  \item Ange hur många olika värden ett tal med $n$ bitar kan anta, och vilket det största är.
  \item Addera två binära tal, och avgöra om summan ryms i åtta bitar.
  \item Beskriva BCD, Graykod, ASCII och 7-segmentkod, och koda och avkoda tal och tecken i dem.
  \item Beräkna en jämn paritetsbit, och förklara vilka fel den upptäcker.
\end{itemize}

\needspace{6\baselineskip}
\subsection*{Frågor att testa dig själv med}
\begin{enumerate}
  \item Varför har den minst signifikanta positionen vikten 1 i \emph{alla} talsystem?
  \item Hur ser du direkt på ett binärt tal om det är jämnt eller udda?
  \item Varför används hexadecimala tal till att skriva binära tal, och inte decimala?
  \item Talet \code{0x100} är ett större tal än \code{0xFF}, men bara med ett. Hur många bitar
    behövs för var och ett av dem?
  \item Vad händer med ett binärt tal när du lägger till en nolla längst till höger? Jämför med
    vad som händer med ett decimalt tal.
  \item Tecknet \code{'7'} och talet 7 är olika saker i en dator. Vad skiljer dem åt, och hur gör
    du om det ena till det andra?
  \item I vilken kod ändras exakt en bit mellan två på varandra följande tal, och varför är det
    användbart?
\end{enumerate}

Talsystemen används i resten av boken, och de sitter först när du har räknat många omvandlingar
själv. Räkna hellre för många än för få. Kalkylatorn i Windows, i läget \textbf{Programmerare},
visar samma tal decimalt, binärt och hexadecimalt samtidigt. Använd den för att kontrollera dina
svar, aldrig för att ta fram dem.

\needspace{6\baselineskip}
\subsection*{Nästa kapitel}
Kapitel~\ref{c2:ch} handlar om logiska grindar och boolesk algebra:
\begin{itemize}
  \item Logiska grindar: NOT, AND, OR, NAND, NOR, XOR och XNOR, och deras sanningstabeller.
  \item Boolesk algebra: notationen, räknelagarna och De Morgans lagar.
  \item Från sanningstabell till grindnät, och från grindnät tillbaka till sanningstabell.
  \item Samma logik byggd med strömkontakter på 24~V, inför Labb~1.
\end{itemize}

% Section 1.4, from lectures/L01/appendix/c_exercises.md.
%
% One deliberate difference: exercise 10 a) spells the text IC74 rather than the name of the class
% the course was first given to, so that the book belongs to no one year.

\section{Övningar}
\label{c1:sec:exercises}\label{c1:app:c}

\begin{aside}
\textbf{Så kontrollerar du ditt arbete.} Alla övningar här räknas med papper och penna. Räkna varje
omvandling för hand först, och kontrollera den sedan genom att räkna tillbaka åt andra hållet: ett
binärt tal du har räknat fram ur ett decimalt ska ge samma decimala tal när du summerar vikterna.
Den kontrollen hittar nästan alla fel, och den finns alltid till hands, också på provet.

Lösningsförslagen finns i bilaga~\ref{app:answers}. Där visas också de vanligaste felaktiga
svaren, och varför de är fel.

Varje övning är märkt med sin sort. Övning~\ref{c1:ex:13} är kapitlets
\textbf{kontrolluppgift}: där räknar du för hand och kontrollerar med kalkylatorn i Windows.
\end{aside}

\exercise{Positionssystem}{Förståelse}
\label{c1:ex:1}
\begin{parts}
  \item Talet 4072 är decimalt. Ange för varje siffra dess position, dess vikt och det värde den
    bidrar med.
  \item I ett talsystem med talbasen 5 finns bara siffrorna 0--4. Vilken vikt har position 2 där?
    Vilket decimalt tal är $243_5$?
  \item Förklara med egna ord varför den högra positionen har vikten 1 oavsett talbas.
\end{parts}

\exercise{Från binärt till decimalt}{Räkna för hand}
\label{c1:ex:2}
Omvandla till decimalt. Använd viktraden 128, 64, 32, 16, 8, 4, 2, 1.
\begin{parts}
  \item \code{1101}
  \item \code{1 0000}
  \item \code{0110 0100}
  \item \code{1000 0001}
  \item \code{1111 1111}
  \item \code{1101 0110}
\end{parts}

\exercise{Från decimalt till binärt}{Räkna för hand}
\label{c1:ex:3}
Omvandla till binärt. Använd vikttabellen för a)--c) och division med 2 för d)--f), och
kontrollera varje svar genom att räkna tillbaka.
\begin{parts}
  \item 13
  \item 37
  \item 100
  \item 200
  \item 255
  \item 1000
\end{parts}

\exercise{Binärt och hexadecimalt}{Räkna för hand}
\label{c1:ex:4}
Omvandla från binärt till hexadecimalt:
\begin{parts}
  \item \code{1010 1111}
  \item \code{0011 1100}
  \item \code{1 1110}
  \item \code{11 1110 1000}
\end{parts}
Omvandla från hexadecimalt till binärt:
\begin{parts}[start=5]
  \item \code{0x7F}
  \item \code{0xC3}
  \item \code{0x2A5}
\end{parts}

\exercise{Hexadecimalt och decimalt}{Räkna för hand}
\label{c1:ex:5}
Omvandla från hexadecimalt till decimalt:
\begin{parts}
  \item \code{0x2A}
  \item \code{0xFF}
  \item \code{0x100}
  \item \code{0x3E8}
\end{parts}
Omvandla från decimalt till hexadecimalt, med valfri metod:
\begin{parts}[start=5]
  \item 75
  \item 160
  \item 4095
\end{parts}

\exercise{Bitar och talområden}{Förståelse}
\label{c1:ex:6}
\begin{parts}
  \item Hur många olika värden kan ett tal med 4, 8, 12 respektive 16 bitar anta, och vilket är det
    största?
  \item Hur många bitar behövs minst för att kunna representera 1000 olika värden?
  \item En nivågivare i en tank ska rapportera nivån i hela steg från 0 till 200. Hur många bitar
    behövs? Hur många av kombinationerna blir då oanvända?
  \item Utan att räkna ut värdet: är \code{1011 0111} ett jämnt eller ett udda tal? Är det större
    eller mindre än 128? Motivera.
\end{parts}

\exercise{Binär addition}{Räkna för hand}
\label{c1:ex:7}
Addera de 8-bitars binära talen. Skriv ut minnessiffrorna, och kontrollera varje summa genom att
göra om talen till decimalt.
\begin{parts}
  \item \code{0011 1010 + 0001 0110}
  \item \code{0111 1111 + 0000 0001}
  \item \code{1100 1000 + 0110 0100}
  \item Vilken av summorna i a)--c) ryms inte i åtta bitar? Vad blir kvar i ett 8-bitars register,
    och vad hände med resten?
\end{parts}

\exercise{BCD}{Räkna för hand}
\label{c1:ex:8}
\begin{parts}
  \item Koda talet 47 i BCD. Koda det sedan som ett 8-bitars binärt tal, och jämför.
  \item Koda årtalet 2026 i BCD. Hur många bitar behövs?
  \item Vilket decimalt tal är \code{1001 0011} i BCD? Vilket tal är samma bitmönster som binärt
    tal?
  \item Är \code{0110 1010} ett giltigt BCD-tal? Motivera.
\end{parts}

\exercise{Graykod}{Förståelse}
\label{c1:ex:9}
\begin{parts}
  \item Bygg Graykoden med tre bitar genom att spegla tvåbitarskoden \code{00, 01, 11, 10}, enligt
    \secref{c1:b3}.
  \item Hur många bitar ändras när en binär räknare går från 7 till 0 (\code{111} till
    \code{000})? Hur många ändras mellan motsvarande Graykoder?
  \item En vinkelgivare med tre bitar står mellan läge 1 och läge 2. Förklara vad styrsystemet kan
    läsa under övergången om givaren använder vanlig binärkod, och vad det kan läsa om den
    använder Graykod.
\end{parts}

\exercise{ASCII}{Räkna för hand}
\label{c1:ex:10}
Använd tabellen i \secref{c1:b4}.
\begin{parts}
  \item Skriv texten \code{IC74} som ASCII-koder, hexadecimalt.
  \item Vilken text är ASCII-koderna \code{0x48 0x65 0x6A}? (\code{e} ligger på \code{0x65}, och
    bokstäverna ligger i ordning.)
  \item Hur mycket skiljer \code{a} och \code{A}? Vilken enda bit skiljer dem åt?
  \item Ett program har läst in siffertecknet \code{7} från ett tangentbord. Hur får programmet
    fram talet 7 ur det? Och hur gör det om talet 3 till tecknet \code{3}?
\end{parts}

\exercise{7-segmentkod}{Konstruktion}
\label{c1:ex:11}
Använd bitordningen g f e d c b a från \secref{c1:b5}.
\begin{parts}
  \item Vilka segment ska lysa för den hexadecimala siffran \code{E}? Ange koden binärt och
    hexadecimalt.
  \item En display får koden \code{0x6D}. Vilken siffra visas?
  \item Segment g har gått sönder och lyser aldrig. Vilka siffror går då inte längre att skilja från
    någon annan siffra? Vilka ser bara konstiga ut?
\end{parts}

\exercise{Paritetsbit}{Räkna för hand}
\label{c1:ex:12}
Använd jämn paritet: paritetsbiten väljs så att det totala antalet ettor blir jämnt, och den skickas
som bit~7 framför de sju ASCII-bitarna.
\begin{parts}
  \item Ange den byte som skickas för tecknen \code{B}, \code{C} och \code{G}.
  \item En mottagare får byten \code{1100 0001}. Har ett fel inträffat? Kan mottagaren veta vilket
    tecken som skulle ha skickats?
  \item Visa med ett eget exempel att två felaktiga bitar i samma byte inte upptäcks.
\end{parts}

\exercise{Kontroll: för hand, och sedan med kalkylatorn}{Kontroll}
\label{c1:ex:13}
\emph{Räkna för hand, kontrollera med ett verktyg, och förklara det du ser.}

Gör delarna i ordning, och skriv ned varje svar innan du går vidare. Poängen går förlorad om du
tittar i kalkylatorn först.

\task{a) För hand}
Omvandla 173 till binärt och hexadecimalt, \code{0x3E8} till decimalt och binärt, och
\code{0b1100 1010} till decimalt och hexadecimalt.

\task{b) Med kalkylatorn}
Öppna kalkylatorn i Windows och välj läget \textbf{Programmerare}. Skriv in vart och ett av talen i
det talsystem det är givet i (klicka först på \code{DEC}, \code{HEX} eller \code{BIN}), och läs av
de andra två. Stämmer dina svar? Om inte: hitta felet i din uträkning, inte bara rätt svar.

\task{c) Något du inte kan förklara än}
Välj \code{DEC}, skriv \code{-1} (tryck \code{1} och sedan \code{+/-}), och läs av \code{HEX} och
\code{BIN}. Byt sedan ordlängd med knappen där det står \code{QWORD} tills det står \code{BYTE}, och
läs av igen. Skriv ned vad du ser i båda fallen.

\task{d)}
Vilket positivt tal har samma bitmönster som det kalkylatorn visade för \code{-1} i läget
\code{BYTE}? Gissa varför datorn skriver \code{-1} just så. Svaret, som heter
\textbf{tvåkomplement}, kommer i kapitel~\ref{c8:ch}.

\exercise{Att lägga till en nolla}{Förståelse, Fördjupning}
\label{c1:ex:14}
\begin{parts}
  \item \code{0b1011} är 11. Vilket tal är \code{0b10110}, där en nolla har lagts till längst till
    höger? Och \code{0b101100}?
  \item Formulera en regel: vad händer med värdet på ett binärt tal när alla bitar flyttas ett steg
    åt vänster och en nolla fylls på från höger?
  \item Vad händer i stället med värdet om den högra biten tas bort och alla bitar flyttas ett steg
    åt höger? Pröva med \code{0b1011} och \code{0b1010}.
  \item Vad motsvarar det att lägga till en nolla längst till höger i ett hexadecimalt tal? Pröva
    med \code{0x2A}.
\end{parts}
Mikrodatorn i boken har instruktioner som gör exakt det här med ett register, och du kommer att
använda dem i kapitel~\ref{c8:ch}.

